\begin{center}
  {\heiti\zihao{2}\PaperTitle\par}
\end{center}
\vspace{0.6em}

\noindent\textbf{摘要：}
本文面向考古工作者在玻璃样品受风化、成分比例失真且检测点存在未检出成分时，仍需获得可解释、可复核分类结论的实际需求，建立“数据清洗—成分变换—风化校正—类型鉴别—亚类划分—关联比较”的完整分析链。原始数据包含58件已分类文物、69个已分类采样点和8件未知样品；按成分和处于85\%--105\%的题设标准剔除15、17号记录后，保留67个有效采样点。对成分数据依次实施闭合、乘法近零替换与CLR变换，并以文物编号分组验证，避免同一文物多点采样造成信息泄漏。

\textbf{针对问题一，}在文物层采用Pearson卡方检验、Monte Carlo精确检验和偏倚校正Cram\'er\'s $V$分析类型、纹饰、颜色与整体表面风化的关系。结果显示，铅钡玻璃与高钾玻璃的风化率分别为70\%和33\%，玻璃类型与风化显著相关（$p=0.0087$，$V_{\mathrm{adj}}=0.321$），而纹饰、颜色尚无充分关联证据。在采样点层，先检验共享精度Dirichlet候选模型；因$\hat R$为1.020--4.896、未通过收敛准入，改用更适合小样本的CLR稳健中心逆扰动，对32个风化点给出保守的风化前校正值。两件整件留出案例的Aitchison距离较同类型中心基准分别降低36.0\%和12.3\%，说明该方法能够保留部分个体组成信息，但不将其夸大为历史成分的精确复原。

\textbf{针对问题二，}以信息增益决策树给出可审计规则：$\mathrm{clr}(\mathrm{PbO})<1.7497$判为高钾玻璃，否则判为铅钡玻璃；再用1个潜变量的PLS--DA复核多成分联合判别，按文物分组交叉验证的准确率、平衡准确率和宏平均F1分别为0.9851、0.9722和0.9807，PbO、BaO、K$_2$O为最重要的助熔相关指标。组内亚类采用MAD筛选5个高变异CLR坐标后实施多起点K-means：高钾玻璃划为4类，平均轮廓系数0.660；铅钡玻璃划为5类，平均轮廓系数0.327。小幅扰动下两类聚类的平均ARI分别为1.000和0.998，同时保留铅钡亚类对特征选择较敏感的结论边界。

\textbf{针对问题三，}将14维CLR--PLS--DA作为主模型，以四变量PLS--DA和PbO决策树作交叉核对。三种模型对A1--A8的结论完全一致：A1、A6、A7为高钾玻璃，A2、A3、A4、A5、A8为铅钡玻璃。对每件未知样品进行100次保持CLR零和约束的局部扰动，并在$\delta=0.01\%,0.02\%,0.04\%$三种近零参数下重算，8件样品均保持100\%原判别；其中A5距边界最近，但最小得分间隔仍为0.221。

\textbf{针对问题四，}在两类玻璃的14维CLR空间分别计算Pearson相关矩阵，以Spearman相关和按文物完整删除检验稳健性。两类均不存在绝对相关系数不小于0.8的成分对；Fe$_2$O$_3$--CuO的类型差异最大，高钾组为0.557、铅钡组为$-0.497$，绝对差为1.054。91组对应相关系数的配对Wilcoxon检验为$p=0.8633$，未发现整体结构存在系统性偏移，因此将少数差异较大的成分对作为配方与风化机制的筛查线索，而不作因果解释。本文最终形成一套从数据质控到未知样品鉴别、从单点校正到结果边界说明的决策流程，可为古代玻璃文物的检测复核、类别初筛和后续考古研究提供定量依据。

\vspace{0.6em}
\noindent\textbf{关键词：}古代玻璃；成分数据；PLS--DA；K-means聚类；敏感性分析

% 摘要（含标题和关键词）不得超过一页。
